% §VII --- Discussion (~0.5 page, ~500 words). Drafted D11.
\section{Discussion}\label{sec:disc}

The empirical results in Section~\ref{sec:exp} call out three
mechanisms behind \method's consistent advantage. First, the
factorized hybrid policy from Proposition~\ref{prop:grad} keeps the
discrete and continuous heads' gradients independent, so the
piecewise reward~\eqref{eq:reward} can simultaneously sharpen the
discrete symbol / channel selection and steer the continuous
trajectory / power without either head crowding the other. The
ablation in Table~\ref{tab:ablation} confirms that switching from the
piecewise reward to a flat $-O(t)$ shaping inflates the joint cost from
$0.30$ to $0.41$ and drops the success rate from $0.81$ to $0.52$
(roughly a one-third relative drop), the largest single regression
among the four ablation axes.

Second, centralizing the critic over the global state mitigates the
non-stationarity that arises when several UAV agents update their
policies concurrently in the same channel. The MAPPO-hybrid baseline
shares this centralized critic but compresses the actor's discrete
heads, so its asymptotic reward in Fig.~\ref{fig:conv} sits roughly
$25\%$ below \method's at convergence; the gap widens as the topology
grows in Fig.~\ref{fig:scale}. The role-aware specialized heads in
\method therefore extract value that the canonical MAPPO weight-sharing
scheme leaves on the table when the agents have heterogeneous decision
types.

Third, the piecewise reward turns the LUT-encoded semantic-feasibility
constraint into a differentiable gradient signal: the third branch
in~\eqref{eq:reward} receives a $+0.5$ trajectory bonus whenever every
UE's similarity exceeds the threshold $\xi_{\mathrm{th}}$, while the
fourth branch penalizes the fraction of UEs that fail. The per-task
similarity CDF in Fig.~\ref{fig:cdf} shows that this design
concentrates probability mass above $\xi_{\mathrm{th}}=0.5$ in a way
that the conventional cost-only baselines (PADDPG, PDQN) cannot match
even at convergence.

The robustness panels (Fig.~\ref{fig:robust}, Fig.~\ref{fig:jam}) and
the large-scale UE generalization in Table~\ref{tab:largeue} reveal a
further property: \method's role-aware actor heads transfer to
unseen channel and topology conditions better than the parameter-shared
variant. We attribute this to the per-head specialization: the four
role groups (UAV symbol, UE symbol, UAV channel, UE channel) absorb
their respective domain shifts independently, so no single head's
distribution drift contaminates the others. The MAPPO-hybrid baseline,
by contrast, propagates a per-agent index embedding through one shared
head, and a perturbation that drifts the input distribution affects
every agent simultaneously.

Two practical implications follow. First, when heterogeneous agents
share a hybrid action space, role-aware specialization is preferable to
parameter sharing despite the parameter count gap; the wall-clock cost
in Table~\ref{tab:wallclock} is modest ($\sim 3\%$ training time
overhead). Second, a piecewise reward grounded in the LUT-encoded
task-fidelity constraint is more effective than a smooth cost-only
shaping, because it keeps the policy gradient informative even when
the cost surface is locally flat or noisy. We expect both insights to
transfer to other heterogeneous-agent semantic-resource scheduling
problems beyond UAV-VQA.
